\subsection{The three exceptional five-point block counts}

\begin{theorem}[Exact affine-five exceptions]\label{thm:five-affine-exceptions}
Every action with exact primitive affine local degree five and
$s\in\{2,3,4\}$ has a complete exponentially small forward row
and a quadratic-exponentially small scalar.
\end{theorem}
\begin{proof}
Write $M=\F_5^s$, $E=U\cap M$, $R=U/E\leq C_4\wr T$.
Since $s\leq4$, the latter group has order prime to five.
Transitivity forces $E\ne1$, since otherwise $5\nmid|U|$.
Schur--Zassenhaus first splits $U=E:R$, then conjugates its
complement into the actual monomial complement in $M:R$.
This does not split $R\to T$. The semisimple monomial modules
$M,M^*$ are cyclic: the orbit of one coordinate vector spans the
module. Every section $W=E/I$ and its dual are therefore cyclic
and can be realized as actual submodules of $M$.

Let $D=\ker(R\to T)$ and $K=\ker(R\to\mathrm{GL}(W))$ for
$W\ne0$. Every coordinate projection of this realized $W$ is
nonzero by transitivity. A diagonal element trivial on $W$ is
therefore one, so $K\cap D=1$ and $K$ embeds normally in $T$.
Its coordinate orbits have one common size $q\mid s$; a fixed
vector on a monomial orbit is determined by one coordinate, giving
$\dim W\leq s/q$. Thus $K=1$ if $\dim W\geq2$ at $s=2,3$,
or if $\dim W\geq3$ at $s=4$. The only further possibility is
$s=4,\dim W=2,q=2$. Then $K\leq C_2^2$ and $T\leq D_8$,
so $R$ is binary. The normal subgroup $K$ has a series of at most
two central $C_2$ factors in $R$: choose an involution in
$K\cap Z(R)$ and then use centrality of a normal order-two group
in the quotient. The same holds after quotienting any normal
$K_0\leq K$.

We need a joint-source onto estimate. If $X=W:S$, with $S$ a
faithful $5'$-group on $W=\F_5^d$ and $W^*$ cyclic, then
\begin{equation}\label{eq:five-cyclic-onto}
 |\Epi(J,X)|\leq5^d|\mathrm{GL}_d(5)|5^{b/5}
                 \qquad(J\leq S_b).
\end{equation}
Indeed fix the single characteristic source subgroup
$F=O^{5'}(J)$ and $B=J/F$. Residuals commute with epimorphisms,
so every map onto $X$ sends $F$ onto $W$ and induces a map
$B\to S$. Put $A=H_1(F,\F_5)$. A surjective equivariant
restriction dualizes to an injection $W^*\hookrightarrow A^*$.
The image of one fixed cyclic vector determines its whole cyclic
$B$-submodule. For each image there are at most
$|\mathrm{GL}_d(5)|$ identifications; faithfulness determines the
top map from each identification. Thus this bounds all top maps
and restrictions jointly by $|\mathrm{GL}_d(5)||A^*|$.
The remaining lift fibre is empty or a torsor for $Z^1(B,W)$;
coprime averaging gives $|Z^1(B,W)|\leq|W|$. Finally
$d_5(F)\leq b/5$ in its actual permutation action, proving
\eqref{eq:five-cyclic-onto}.
A central $C_2$ lift has at most $|\Hom(J,C_2)|\leq2^{b/2}$
choices over a fixed map, whether split or not. The two central
steps above therefore cost at most $2^b$ in total.

Fix now $N\lhd U$ and $I=N\cap E$. If $I=E$, these terms sum
exactly to $Z_J(R)$. If $\dim(E/I)=1$, retain the whole comparator
$U/I$. For this line let $\chi:R\to\F_5^*$ be its action.
The map $(\chi,\pi_T)$ is injective since $K\cap D=1$, so
$U/I\hookrightarrow\operatorname{AGL}_1(5)\times T$ has a
faithful action of degree $5+s\leq4s$.
For $\dim(E/I)\geq2$, the normal $N/I$ is disjoint from $W$
and centralizes it. It lies in $W\times K$ and projects injectively
to the $5'$-group $K$; its graph into $W$ is trivial. Hence
$N/I=K_0\lhd R$. Apply \eqref{eq:five-cyclic-onto} if $K=1$,
or apply it after the at most two central binary steps otherwise.

This proves, with a fixed list $\mathcal Q_U$ consisting of $R$
and all literal line quotients $U/I$, that
\[
 Z_J(U)\leq\sum_{Q\in\mathcal Q_U}Z_J(Q)+C_U2^{\eta_sb},
 \qquad C_U=\nu(U)5^s|\mathrm{GL}_s(5)|,
\]
where $\eta_2=\eta_3=\log5/5$ and
$\eta_4=1+\log5/5$. Equal abstract quotients at different $I$
remain different entries of this fixed list. Pad every comparator
to faithful degree $v=4s$.

For clarity the finite comparator argument extends directly to
this fixed list. Let $m=|\mathcal Q_U|$ and
$S_U(J)=m^{-1}\sum_{Q\in\mathcal Q_U}Z_J(Q)$. Jensen and
Lemma~\ref{lem:quotient-moment} give, for every integer $q\geq1$,
\[
 \sum_{J\leq S_b}S_U(J)^q
 \leq m^{-1}\sum_Q\sum_J Z_J(Q)^q\leq s_{b+qv}.
\]
The proof of Theorem~\ref{thm:finite-comparator} uses precisely
this moment inequality and positivity, so applies with statistic
$S_U$ and multiplier $m$. Here $h(5s)=10,14,20$ and
$v=8,12,16$, respectively; $v<h(5s)$ and
$\eta_s<h(5s)/8$, since $\log5/5<7/15$.
The finite union therefore has the claimed bounds with its original
physical weights and one complete residual terminal.
\end{proof}
